\section{Four distinct blocks: an exact certificate proof}\label{sec:reach-four}

The next reach bound has the same uniform-input constant, but its proof
uses a finite family of exact linear identities. We give the complete
necessary-condition relaxation and its symmetry reduction before applying
the certificates. The computer-assisted part establishes only the weighted
pair inequality below; the passage from that inequality to a four-block
schedule is analytic. No time discretization of the input is involved.

\begin{theorem}[Four distinct blocks, computer-assisted]\label{thm:reach-four}
For $n\ge5$, four distinct modes suffice to reach $\min\{T,B_4(n,E)\}$
with one-sided error at most $E$.
\end{theorem}

\subsection{The weighted pair inequality}

Use the reach maps~\eqref{eq:reach-definition} on $[0,L]$. Let
$R_i=\Phi_i(0)$ and define the pair maxima
\begin{equation}\label{eq:pair-exclusions-four}
 \begin{split}
 M&=\max_{a\ne b}\Phi_b(R_a),\\
 P_i&=\max_{a,b\notin\{i\},\ a\ne b}\Phi_b(R_a),\\
 Q_{ij}=Q_{ji}&=
 \max_{a,b\notin\{i,j\},\ a\ne b}\Phi_b(R_a),\qquad i\ne j.
 \end{split}
\end{equation}
These sets of pairs are nonempty. Assume $M<L$, so all pair and first
reaches are uncapped. For a three-element subset $S\subset[n]$, write
\begin{equation}\label{eq:H-S}
 \mathcal H_S=\sum_{i\in S}\left(P_i+\sum_{j\ne i}Q_{ij}\right).
\end{equation}
An excluded pair wholly in $S$ is counted twice in this sum.

\begin{lemma}[Weighted pair inequality, computer-assisted]\label{lem:weighted-pair}
For $n\ge5$, $M<L$, any three-element set $S$, and any distinguished
mode $z$,
\begin{equation}\label{eq:weighted-pair-strong}
 \mathcal H_S\ge \frac{3n^2}{n-1}E
       +\frac{3n^2}{(n-1)^2}\sum_{i\ne z}R_i.
\end{equation}
In particular,
\begin{equation}\label{eq:weighted-pair}
 \mathcal H_S\ge3nB_2(n,E).
\end{equation}
\end{lemma}

To see the implication between the two displayed inequalities, choose $z$
with smallest first reach. Then
\[
 R_z=\sum_i A_i(R_z)\le\sum_i A_i(R_i)
     =\sum_i(R_i-E),
 \qquad \sum_{i\ne z}R_i\ge nE.
\]
Substitution in~\eqref{eq:weighted-pair-strong} gives
\eqref{eq:weighted-pair}, since
$B_2=nE(2n-1)/(n-1)^2$. The stronger statement is what the certificates
prove, and does not assume that $z$ is smallest.

\subsection{The complete necessary-condition linear program}\label{subsec:event-lp}

Scale time by $E$, so that $E=1$. For this subsection and its implementation
only, label modes $0,\ldots,n-1$ and set $S=\{0,1,2\}$. The variables are
nonnegative $R_i$, and event times $t_v$ and allocations $a_{v,i}$ for
\[
 v\in\mathcal V=
 \{M\}\ \cup\ \{P_i:i\in S\}\ \cup\
 \{Q_{ij}:i<j,\ \{i,j\}\cap S\ne\varnothing\}.
\]
There are $1+3+[3(n-3)+3]=3n-2$ events. For event $v$, let $D_v$ be
its excluded set: $\varnothing$, $\{i\}$, or $\{i,j\}$ as appropriate.
Choose an unordered pair $P=\{p,q\}$ attaining the global pair maximum
$M$. The relaxation consists of exactly the following constraints:
\begin{align}
 &\sum_{i=0}^{n-1}a_{v,i}=t_v &&(v\in\mathcal V),
       \label{eq:event-mass}\\
 &R_j-t_v+a_{v,k}\le-1
       &&(v\in\mathcal V;\ j,k\notin D_v,\ j\ne k),
       \label{eq:event-pair}\\
 &a_{Q_{ij},k}\le a_{P_i,k}
       &&(i\in S,\ j\ne i;\ 0\le k<n),
       \label{eq:event-inclusion}\\
 &a_{v,k}\le a_{M,k}
       &&(v\in\mathcal V\setminus\{M\};\ 0\le k<n),
       \label{eq:event-global}\\
 &a_{M,k}\le a_{Q_{ij},k}
       &&(\{i,j\}\cap P=\varnothing,\ Q_{ij}\in\mathcal V;\ 0\le k<n),
       \label{eq:event-equal-Q}\\
 &a_{M,k}\le a_{P_i,k}
       &&(i\in S\setminus P;\ 0\le k<n).
       \label{eq:event-equal-P}
\end{align}
All variable nonnegativity constraints are imposed as well. The objective is
\begin{equation}\label{eq:event-objective}
 C=(n-1)^2\left(\sum_{i\in S}t_{P_i}
   +\sum_{\substack{i<j\\\{i,j\}\cap S\ne\varnothing}}
          |\{i,j\}\cap S|\,t_{Q_{ij}}\right)
          -3n^2\sum_{i\ne z}R_i.
\end{equation}
We will certify $C\ge3n^2(n-1)$.

Every input in the uncapped case gives a feasible point: set $t_v$ to its
pair reach and $a_{v,i}=A_i(t_v)$. Equation~\eqref{eq:event-mass}
expresses total allocation. For~\eqref{eq:event-pair}, the event reach
is at least $\Phi_k(R_j)$ and this endpoint is uncapped. Hence
$H_k(t_v)\ge R_j+1$. Excluding more modes cannot increase a pair maximum,
which gives~\eqref{eq:event-inclusion} and~\eqref{eq:event-global} by
allocation monotonicity. If an exclusion avoids the maximizing pair $P$,
that pair remains available and its event time equals $M$. The reverse
allocation inequalities~\eqref{eq:event-equal-Q}--\eqref{eq:event-equal-P}
express these equalities. Together with~\eqref{eq:event-mass}, the two
allocation orders also force the event times to be equal.

There are no additional constraints: in particular we impose neither
first-root equations nor an ordering of the $R_i$, and no chronological
ordering between unrelated events. The relaxation can therefore include
nonphysical points. A lower bound valid on this larger feasible set is
nevertheless valid for every control. No converse representation claim
is used.

\subsection{Exhaustive symmetry reduction}\label{subsec:pair-symmetry}

Permutations preserving $S$ carry a maximizing pair to one of three
representatives, according to its intersection size with $S$. Permutations
preserving both $S$ and this pair then place $z$ in the representatives
shown in Table~\ref{tab:pair-types}. For example, when $P=\{0,3\}$,
the four possible positions of $z$ are the singleton $S\cap P$, one of
the two elements of $S\setminus P$, the singleton $P\setminus S$, and
an element outside $S\cup P$. These are exactly the four displayed cases.
The same partition argument proves the other two rows. Ties between
maximizing pairs do not matter: any maximizing pair may be chosen.

\begin{table}[tb]
\centering
\begin{tabular}{lll}
\toprule
$P$ & Representatives for $z$ & Range\\
\midrule
$\{0,1\}$ & $0,2,3$ & $n\ge5$\\
$\{0,3\}$ & $0,1,3,4$ & $n\ge5$\\
$\{3,4\}$ & $0,3$ & $n\ge5$\\
$\{3,4\}$ & $5$ & $n\ge6$\\
\bottomrule
\end{tabular}
\caption{All ten types of $(S,P,z)$ after relabeling $S=\{0,1,2\}$.}
\label{tab:pair-types}
\end{table}

For one type, keep the indices in $K=S\cup P\cup\{z\}$ individually
labeled and let the remaining indices be interchangeable. Averaging a
feasible point over all permutations of those remaining indices preserves
every constraint and the objective. Thus restricting variables to be
constant on permutation orbits preserves the infimum of the relaxation.
This assertion needs no optimizer or boundedness assumption: every feasible
point has an invariant feasible average with the same value.

Here is an explicit orbit rule sufficient to reconstruct the quotient.
A variable is identified by $(R,i)$, $(t,v)$, or $(a,v,i)$. Retain each
index in $K$ verbatim. Replace each other index, in order of first occurrence
in this identifier, by $b_0,b_1,\ldots$, reusing its symbol on subsequent
occurrences. A $Q$ event contains at most one interchangeable index,
because its excluded pair meets $S$. Thus this rule preserves exactly
the equalities and inequalities of the unspecified indices, without an
orientation ambiguity for its unordered pair. In particular, allocation
to the interchangeable index excluded by a $Q$ event and allocation to a
different interchangeable index remain separate variables.

Replace variables in each row of
\eqref{eq:event-mass}--\eqref{eq:event-equal-P} by these orbit labels,
add coefficients that now refer to the same variable, and delete duplicate
rows. In the objective, add all multiplicities of the same variable orbit.
This defines the quotient matrices and objective uniquely, independently
of any numerical solver. The bundled checker implements precisely this
construction. Its ordering enumerates $R_i$ first, then $M$, the $Q_{ij}$
in lexicographic order, and the three $P_i$; it retains the first occurrence
of each variable orbit and each distinct row. The certificate coordinates
refer to that ordering.

At most six indices lie in $K$. A pair-constraint row involves at most
three interchangeable indices: one from its event and its two available
modes. An allocation-order row involves at most two. Thus all variable and
row patterns are already present once $n\ge9$. Additional interchangeable
indices repeat those patterns. They affect only multiplicities in mass
equations and the objective. If $a=|K|$, a mass row has $n-a$ identical
allocation terms, or one separately labeled term and $n-a-1$ others when
an interchangeable index is excluded by the event. Hence its coefficients
are affine functions of $n$. Objective orbit multiplicities are constant
or $n-a$, multiplied by $(n-1)^2$ for event times or by $-3n^2$ for
first reaches. All objective coefficients are consequently integer
polynomials of degree at most three. The inequality matrix and right-hand
sides are constant for $n\ge9$.

For $n\ge9$, the quotient has between 57 and 162 variables, 151 and 826
inequality rows, and 10 and 19 equality rows, depending on the type. For the symbolic
check, the exact affine mass coefficients and exact affine objective
multiplicities are reconstructed from the quotient at $n=9,10$. This is
reconstruction of the established affine counts, not interpolation of an
unknown function. All individually labeled indices are at most 5, and
three interchangeable indices suffice for every row pattern. First-occurrence
ordering of patterns is therefore stable when further indices are added.

\subsection{Exact dual identities and the verification boundary}\label{subsec:pair-certificates}

Write the quotient as $Ax\le b$, $B(n)x=d$, $x\ge0$, with objective
$c(n)^Tx=C$. In this notation $b$ has entries zero or minus one and $d=0$.
For a finite dimension, a certificate supplies a positive integer
denominator $D$ and integer vectors $Y\le0$, $Z$ satisfying
\begin{equation}\label{eq:finite-dual}
 D c=A^TY+B^TZ,\qquad b^TY+d^TZ=3n^2(n-1)D.
\end{equation}
For every feasible $x$, multiplying its inequalities by the nonpositive
entries of $Y$ reverses their direction. Therefore
\[
 D C=Y^TAx+Z^TBx\ge Y^Tb+Z^Td=3n^2(n-1)D.
\]
No dual optimality assertion is required. The certificates have zero
coefficient residual, so nonnegative-variable residual multipliers are
unnecessary.

For $5\le n\le22$, the data contain all nine types at $n=5$ and all
ten at each $n=6,\ldots,22$: a total of $9+17\cdot10=179$ finite
certificates. The checker verifies the case identifiers, signs, integrality,
every coordinate of~\eqref{eq:finite-dual}, and the stated lower bound.

For each type, a further certificate gives integer polynomials
$D(n)$, $D_0(n)$, $Y(n)$, and $Z(n)$ satisfying
\begin{align}
 D(n)&=(n-1)^2D_0(n),\label{eq:poly-D}\\
 D_0(n)c(n)&=A^TY(n)+B(n)^TZ(n),\label{eq:poly-coordinate}\\
 b^TY(n)+d^TZ(n)&=3n^2(n-1)D_0(n).\label{eq:poly-bound}
\end{align}
The checker verifies these coefficient by coefficient, using integer
polynomial addition and multiplication. It expands $D(23+t)$ and each
entry of $-Y(23+t)$. The former has positive constant coefficient and
nonnegative remaining coefficients; the latter have only nonnegative
coefficients. Consequently $D(n)>0$ and $Y(n)\le0$ for every real
$n\ge23$, and~\eqref{eq:poly-D} gives $D_0(n)>0$ there. Dividing
\eqref{eq:poly-coordinate}--\eqref{eq:poly-bound} by $D_0(n)$ gives
the same dual bound for every integer $n\ge23$. These checks cover
all dimensions $n\ge5$ and prove~\eqref{eq:weighted-pair-strong}.

The supplement's \path{verification/reference/} directory contains the
complete data in \path{certificates_general_four_block.json} and the
checker \path{verify_general_four_block.py}. Run, from the paper directory,
\begin{quote}\small\ttfamily
python verification/reference/verify\_general\_four\_block.py
\end{quote}
The checker requires only the Python standard library and uses integer
arithmetic throughout. It reads certificate arrays as data, not executable
expressions. Numerical optimization and symbolic algebra were used to
discover the certificates; neither a numerical solution, a floating-point
tolerance, nor an optimization status is trusted by the checker. Its
successful output verifies all 179 finite cases and all ten polynomial
cases. A provenance manifest records SHA-256 hashes of the bundled
original artifacts. The proof's computational dependency is this finite
data and its exact verification, together with the mathematical reduction
and orbit coverage above. It is explicitly a computer-assisted proof.

\subsection{From weighted pairs to four distinct blocks}

\begin{proof}[Proof of Theorem~\ref{thm:reach-four}]
Put $L=\min\{T,B_4\}$ and assume no sequence of at most four distinct
modes reaches $L$. All first, pair, and triple reaches are then uncapped.
Let $G$ be the maximum three-distinct-mode reach and $N_i$ the corresponding
maximum excluding $i$. Theorem~\ref{thm:reach-three} implies $G\ge B_3$:
if $L\le B_3$, that theorem already contradicts the assumed failure.

For each $k\ne i$, appending $k$ to a pair attaining $Q_{ik}$ gives
a triple excluding $i$ and ending no later than $N_i$. Hence
$A_k(N_i)\le N_i-E-Q_{ik}$. Summing over $k\ne i$ gives
\[
 A_i(N_i)\ge(n-1)E+\sum_{k\ne i}Q_{ik}-(n-2)N_i.
\]
Appending $i$ to a pair attaining $P_i$ similarly gives
$A_i(G)\le G-E-P_i$. As $N_i\le G$, allocation monotonicity implies
\begin{equation}\label{eq:excluded-triple-ineq}
 (n-2)N_i\ge nE+P_i+\sum_{k\ne i}Q_{ik}-G.
\end{equation}
Choose a triple attaining $G$ and let $S$ be its mode set. Then $N_i=G$
outside $S$. Sum~\eqref{eq:excluded-triple-ineq} over its three elements:
\[
 \sum_iN_i\ge\left(n-3-\frac3{n-2}\right)G
                  +\frac{3nE+\mathcal H_S}{n-2}.
\]
The coefficient of $G$ is positive for $n\ge5$. Applying
$G\ge B_3$ and Lemma~\ref{lem:weighted-pair} yields
\begin{equation}\label{eq:aggregate-triple}
 \sum_iN_i\ge nB_3,
\end{equation}
because $n(B_2+E)=(n-1)B_3$.

Finally, appending $i$ to a maximizing triple excluding it uses four
distinct modes. Its failure implies $H_i(L)>N_i+E$. Summing gives
\[
 (n-1)L>\sum_iN_i+nE\ge nB_3+nE=(n-1)B_4,
\]
contradicting $L\le B_4$. The strict inequality follows from latest
endpoints even when $H_i$ has flat portions. Endpoint monotonicity makes
the endpoint constraints control the entire one-sided error trajectory;
unselected modes have $W_i-A_i=-A_i\le0$.
\end{proof}

An earlier five-mode proof used a separate weighted-pair event relaxation.
Its verifier and certificate data are retained in the supplement as an
independent check of the smallest case. The all-dimension certificates
above subsume that positive result. Neither the five-mode checker nor the
thirty-case three-block checker is needed to establish the theorem as
presented here.
